\documentclass[12pt,a4paper]{article}
\usepackage{ugmath}
\usepackage{float}
\usepackage{placeins}
\usepackage[labelfont=bf,labelsep=space]{caption}
\newcommand{\paren}[1]{\left( #1 \right)}
\newcommand{\alumno}{Ricardo León Martínez}
\newcommand{\materia}{Fisica I}
\newcommand{\profesor}{Lucero Uscanga Aguilera}
\newcommand{\tarea}{Tarea 5}
\newcommand{\fecha}{13/3/2026}

\begin{document}
    \begin{center}
    {\large\textbf{UNIVERSIDAD DE GUANAJUATO}}\\[0.3cm]
    {\normalsize\textbf{DIVISIÓN DE CIENCIAS NATURALES Y EXACTAS}}\\
    {\normalsize\textbf{CAMPUS GUANAJUATO}}\\[1cm]

    {\Large\textbf{\tarea\ (\materia)}}\\[1cm]
    \end{center}

    \textbf{Nombre:} \alumno \hfill 
    \textbf{Fecha:} \fecha \hfill 
    \textbf{Calificación:} \rule{3cm}{0.4pt} \\[0.3cm]

    \begin{problema}
        Una masa de 10kg se mueve bajo la acción de la fuerza $\vec{F}=\hat{u}_{x}(5t)+\hat{u}_{y}(3t^{2}-1)$N.
        Cuando $t=0$ el cuerpo está en reposo en el origen.
        \begin{enumerate}
            \item[(a)] Hallar el momento y la energía cinética del cuerpo cuando $t=0$s.
            \item[(b)] Calcular el impulso y el trabajo efectuado por la fuerza de $t=0$ a $t=10$s.
            Comparar con las repuestas en (a). 
        \end{enumerate}
    \end{problema}

    \begin{problema}
        Se lanza verticalmente hacia arriba un cuerpo de 20kg de masa con una velocidad de
        50m/s. Calcular
        \begin{enumerate}
            \item[(a)] Los valores iniciales de $E_{k},E_{p}$ y $E$
            \item[(b)] $E_{k}$ y $E_{p}$ despues de 3s
            \item[(c)] $E_{k}$ y $E_{p}$ a 100m de altura
            \item[(d)] La altura del cuerpo cuando $E_{k}$ es reducida un 80\% de su valor inicial.   
        \end{enumerate}
    \end{problema}

\end{document}